\documentclass[]{article}
\usepackage{../../../tex/packages/hebrewSupport}
\usepackage{setspace}
\usepackage[margin=0.8in]{geometry}
\usepackage{multicol}
\setlength{\parindent}{0pt}
\setlength{\parskip}{8pt}



\author{שחר פרץ}
\title{מחויבות אישית $\sim$ י''א $\sim$ רפלקציה הייפר־ממדית}
\begin{document}
	\nonumber
	\onehalfspacing
	
	\maketitle
	
	
	\begin{center}
		\begin{multicols}{4}
			\textbf{שם התלמיד: }שחר פרץ\\
			\textbf{מספר תעודת זהות: }334558962\\
			\textbf{כיתה: }י''א 5\\
			\textbf{שם המחנך: }הדר פרי אגמון
		\end{multicols}
	\end{center}
	
	\large
	
	בממד העיוני, \\
	ימי המיקוד הציגו לי את האופציה הזו, ובכך עזרו לי למצוא את ההתנסות האישית הנכונה עבורי. היה לי את מי לשתף בחוויות שלי, ויכולתי לפנות לעזרה ובקשת מידע למחנכת ולרכז השכבה. 
	
	בממד האישי, \\
	מקום ההתנסות שלי היה להיות אחראי מחויבות אישית בכפר הירוק. במסגרת ההתנסות, הייתי צריך לבדוק שכל תלמיד ממלא את חובותיו במקום מסוים, לענות לשאלות טכניות של תלמידים ולאסוף טפסים ולארגנם בשביל המחנכת (הדר פרי אגמון). חווית הצלחה בעבורי הייתה לראות לקראת סוף השנה את הכל מאורגן – אני יודע איפה כל תלמיד עשה את המחויבות האישית שלו, באיזה תאריך הוא יסיים ויגיש את הטפסים הרלוונטיים, וכו'. התעורר הקושי לקבל תגובה בפרטי מכל אחד, בנוגע למספר עניינים שעסקו במחויבות אישית, אך לבסוף כולם ענו. 
	
	בממד הקבוצתי, \\
	אופי התפקיד היה מוגדר מראש, ולכן לא בחרתי להרים אף פרויקט. פרויקט שהיה עלי להעביר היה איסוף כל החתימות של כל התלמידים תוך כשבוע. קושי שהיה בעבורי הוא שמכיוון ואיסוף החתימות מצריך פניה אל אנשים שלרוב נמצאים מחוץ לכפר והגישה אליהם אינה מהירה, ולכן נדרשו פתרונות יצירתיים על־מנת להשיג את כל הטפסים במסגרת הזמן הנדרשת. 
	
	
	
	\ndoc
	
\end{document}